\begin{frame}{Looking ahead}
  % \framesubtitle{And why the rest of this week is about it}
  \footnotesize
  \setlength{\itemsep}{1pt}
  \setlength{\parsep}{0pt}

  \begin{itemize}
    \item openADR is, structurally, a \textbf{Communication- and Function-layer} technology: messages, roles, processes, systems.
    \begin{itemize}
      \item Agreements there are \emph{protocol-level}: endpoint, format, sequence, security.
    \end{itemize}
    \pause
    \item The DSO sits on top of a large existing information layer: \alert{CIM, CGMES, IEC 61850, grid data models}.
    \begin{itemize} 
      \item openADR assets (chargers, PV, storage) will be \emph{inside} that layer as soon as they matter for grid calculations — and Day 2 starts there.
    \end{itemize}
    \pause
    \item The \textbf{Information layer is where openADR is thinnest}:
    \begin{itemize}
      \item The information standardization in which openADR finds itself is quite nuanced.
      \item In order to achieve interoperability with adjacent systems, openADR will have to adopt that nuance
    \end{itemize}

  \end{itemize}

  \vspace{0.15cm}
  \begin{alertblock}{The week's thesis}
    \emph{openADR tells DER what to do; CIM/CGMES tell everyone what the grid \emph{is}. You need both — and a bridge between them.}
    \vspace{0.1cm}
    \alert{Today's goal:} \textbf{understand what it takes to design an information standard that aligns with adjacent standards.} 
  \end{alertblock}
\end{frame}
